\exercisesection{}

\begin{enumerate}
\def\labelenumi{\arabic{enumi})}
\tightlist
\item
  设 \(\mathbf{T}\) 是一个集合，\(p\) 是 \(\mathbf{T}\) 上的负荷。对于 \(\mathbf{A}\) 的每个子集 \(\mathbf{T}\)，置 \(p(\mathbf{A}) = p(\varphi_{\mathbf{A}})\)。
\end{enumerate}

\begin{enumerate}
\def\labelenumi{\alph{enumi})}
\item
  当 A 包含于 B 时，\(p(A) \leqslant p(B)\)。
\item
  对于 \((\mathbf{A}_n)\) 的任意有限或无限子集序列 \(\mathbf{T}\)，有 \(p(\bigcup \mathbf{A}_n) \leqslant \sum p(\mathbf{A}_n)\)。
\item
  对于 T 的任意递增子集序列 \((\mathrm{A}_n)_{n\geqslant 1}\)，有 \(p\big(\bigcup_{n}\mathrm{A}_{n}\big) = \lim_{n\to \infty}p(\mathrm{A}_n)\)。
\end{enumerate}

\begin{enumerate}
\def\labelenumi{\arabic{enumi})}
\setcounter{enumi}{1}
\tightlist
\item
  设 \(\mathbf{T}\) 是一个集合，\(p\) 是 \(\mathbf{T}\) 上的负荷。如果 \(f \in \mathcal{F}_{+}(\mathbf{T})\)（分别地，如果 A 是 \(\mathbf{T}\) 的子集），且 \(p\)（分别地，\(p(f) = 0\)），则称 f（分别称 A）为 \(p(\mathbf{A}) = 0\)-可忽略的。
\end{enumerate}

\begin{enumerate}
\def\labelenumi{\alph{enumi})}
\item
  设 \(f\) 和 \(g\) 属于 \(\mathcal{F}_{+}(\mathrm{T})\)，且 t 是满足 \(t\)\(f \leqslant t \cdot g\) 的正实数。若 g 是 \(g\)\(p\)\(f\)\(p\)-可忽略的，则 f 也是如此。\(p\)-可忽略函数序列的和与上包络也是 -可忽略的。
\item
  \(p\)-可忽略集的每个子集都是 \(p\)-可忽略的，\(p\)-可忽略集的有限或无限序列的并也是如此。
\item
  设 \(f \in \mathcal{F}_+(\mathrm{T})\)。对于每个有限实数 \(a > 0\)，记 \(\mathbf{T}_a\) 为满足  的 \(t \in \mathrm{T}\) 的集合。证明 \(f(t) \geqslant a\)\(p(\mathbf{T}_a) \leqslant a^{-1} \cdot p(f)\)。由此推出，f 是 \(f\)\(p\)-可忽略的，当且仅当满足  的 \(t \in \mathrm{T}\) 的集合是 \(f(t) \neq 0\)\(p\)-可忽略的。
\item
  设 \(f \in \mathcal{F}_{+}(\mathrm{T})\)。如果 \(p(f)\) 有限，则满足 f(t) 为无穷的 \(t \in \mathrm{T}\) 的集合是 \(f(t)\)\(p\)-可忽略的。
\item
  设 \(f\) 和 \(g\) 属于 \(\mathcal{F}_{+}(\mathrm{T})\)。如果满足 \(t \in \mathrm{T}\) 的 \(f(t) \neq g(t)\) 的集合是 \(p\)-可忽略的，则 \(p(f) = p(g)\)。
\end{enumerate}

\begin{enumerate}
\def\labelenumi{\arabic{enumi})}
\setcounter{enumi}{2}
\item
  设 \(\mathbf{T}\) 是一个集合，\(p\) 是 \(\mathbf{T}\) 上的负荷。对于 \((f_n)_{n\geqslant 1}\) 中元素的每个序列 \(\mathcal{F}_{+}(\mathbf{T})\)，有 \(p(\liminf_{n\to \infty}f_n)\leqslant \liminf_{n\to \infty}p(f_n)\)。
\item
  设 \(\mathbf{T}\) 是一个集合，\(\mathbf{M}\) 是 \(\mathbf{T}\) 上的负荷，\(\mathbf{E}\) 是实或复巴拿赫空间，p 是满足 \(p\)\(p \geqslant 1\) 的有限实数。记 \(\mathcal{F}\) 为从 \(\mathbf{T}\) 到 \(\mathbf{E}\) 的映射集合，并且对于 \(f\)\(\mathcal{F}\) 或 \(\mathcal{F}_{+}(\mathbf{T})\) 中的 f，置 \(\mathrm{N}_p(f) = \mathrm{M}(|f|^p)^{1/p}\)。
\end{enumerate}

\begin{enumerate}
\def\labelenumi{\alph{enumi})}
\item
  对于 \((f_n)_{n\geqslant 1}\) 中元素的每个序列 \(\mathcal{F}_{+}(\mathrm{T})\)，有 \(\mathrm{N}_p\left(\sum_{n\geqslant 1}f_n\right)\leqslant \sum_{n\geqslant 1}\mathrm{N}_p(f_n)\)
\item
  令 \(\mathcal{F}^p\) 为满足 \(f \in \mathcal{F}\) 有限的 \(\mathrm{N}_p(f)\) 的集合。证明 \(\mathcal{F}^p\) 是 \(\mathcal{F}\) 的线性子空间，\(\mathrm{N}_p\) 是 \(\mathcal{F}^p\) 上的半范数。证明，满足 \(\mathcal{N}\) 的 \(f \in \mathcal{F}^p\) 所成的集合 \(\mathrm{N}_p(f) = 0\)，等于在某个 M-可忽略集之外为零的 \(f \in \mathcal{F}\) 所成的集合。证明赋范空间 \(\mathcal{F}^p / \mathcal{N}\) 是完备的（参见第 IV 章，第 3 节，第~3号）。
\end{enumerate}

\begin{enumerate}
\def\labelenumi{\arabic{enumi})}
\setcounter{enumi}{4}
\item
  设 \(\mathbf{T}\) 是一个集合。记 \(\mathcal{B}_{+}(\mathbf{T})\) 为 \(\mathbf{T}\) 上有界正实函数的集合，记 \(p_0\) 为从 \(\mathcal{B}_{+}(\mathbf{T})\) 到  的区间 \([0, +\infty]\) 的映射，它满足第~1号定义 1 的条件 a) 至 d)。对于每个函数 \(\overline{\mathbf{R}}\)\(a)\)\(d)\)\(f \in \mathcal{F}_{+}(\mathbf{T})\)，记 \(p(f)\) 为数集 \(p_0(f_0)\) 的上确界，其中 \(f_0\) 遍历所有满足 \(\leqslant f\) 的有界函数。证明 p 是 \(p\)\(\mathbf{T}\) 上的负荷。将 \(\mathcal{B}_{+}(\mathbf{T})\) 换成 \(\mathbf{T}\) 上（有限）正实函数的集合后，问题同样成立。
\item
  设 \(\mathbf{T}\) 是一个集合，\(\mathcal{I}_{+}\) 是 \(\mathcal{F}_{+}(\mathrm{T})\) 的子集，\(\mathbf{M}\) 是从 \(\mathcal{I}_{+}\) 到 \(\overline{\mathbf{R}}_{+}\) 的映射。作如下假设：
\end{enumerate}

\begin{enumerate}
\def\labelenumi{\alph{enumi})}
\item
  对于每个 \(f \in \mathcal{I}_+\) 和每个实数 \(t \geqslant 0\)，函数 \(t \cdot f\) 属于 \(\mathcal{I}_+\)，且 \(M(t \cdot f) = t \cdot M(f)\)。
\item
  如果 f 和 g 属于 \(f\)\(g\)\(\mathcal{I}_+\)，那么 \(f + g\)、\(\sup(f, g)\) 和 \(\inf(f, g)\) 也属于 ，且有 \(M(f + g) = M(f) + M(g)\)。
\item
  对于 \((f_n)_{n\geqslant 1}\) 中元素的每个递增序列 \(\mathcal{I}_+\)，函数 \(f = \sup f_n\) 属于 \(\mathcal{I}_+\)，且 \(\mathrm{M}(f) = \lim_{n\to \infty}\mathrm{M}(f_n)\)。
\end{enumerate}

对于每个函数 \(f \in \mathcal{F}_{+}(\mathrm{T})\)，置 \(p(f) = \inf_{g \in \mathcal{I}_{+}, g \geqslant f} \mathrm{M}(g)\)。证明 p 是 \(p\)\(\mathrm{T}\) 上的负荷（参见第 IV 章，第 1 节，第~3号）。如果 f 和 g 属于 \(f\)\(g\)\(\mathcal{F}_{+}(\mathrm{T})\)，且 \(\mathrm{A}\) 和 \(\mathrm{B}\) 是 \(\mathrm{T}\) 的子集，则 \(p(\sup(f, g)) + p(\inf(f, g)) \leqslant p(f) + p(g)\)，并且 \(p(\mathrm{A} \cup \mathrm{B}) + p(\mathrm{A} \cap \mathrm{B}) \leqslant p(\mathrm{A}) + p(\mathrm{B})\)。

\begin{enumerate}
\def\labelenumi{\arabic{enumi})}
\setcounter{enumi}{6}
\tightlist
\item
  在 Lindelöf 空间上，每个测度、每个函数和每个子集都是可控的。
\end{enumerate}

¶ 8) 再次使用第 IV 章第 1 节习题 5 的记号。回忆 \(\mu = \sum_{k=-\infty}^{\infty} \sum_{n=1}^{\infty} n^{-3} \varepsilon_{(1/n,k/n^2)}\)，且 \(\mu^{*}(D) = +\infty\)。证明 \(\mu^{\bullet}(D) = 0\)；由此推出 \(\mu\) 不是可控的，尽管它是具有两两不交紧支集的测度级数之和。证明存在将 X 划分为 Borel 子集  的划分 \((X_n)_{n \in \mathbb{N}}\)，使得对于每个 \(X\)，\(X_n\)\(\mu^{\bullet}(X_n)\) 都有限。\(n \in \mathbb{N}\)

\begin{enumerate}
\def\labelenumi{\arabic{enumi})}
\setcounter{enumi}{8}
\item
  设 \(\mathbf{T}\) 是拓扑空间。如果 \(\lambda\) 和 \(\mu\) 是 \(\mathbf{T}\) 上的两个正测度，且对于 \(\lambda^{\bullet}(\mathrm{U}) = \mu^{\bullet}(\mathrm{U})\) 中的每个开集 \(\mathbf{U}\) 都有 \(\mathbf{T}\)，则 \(\lambda = \mu\)（先证明对于 \(\lambda^{\bullet}(f) = \mu^{\bullet}(f)\) 上每个满足 \(f \geqslant 0\) 的下半连续函数 f，都有 \(\mathbf{T}\)，再证明 \(\lambda^{\bullet} = \mu^{\bullet}\)）。
\item
  设 \(\mathbf{T}\) 是拓扑空间，\(\mu\) 是 \(\mathbf{T}\) 上的正测度；假设 \(\mathbf{T}\) 是一列 \(\mu\)-可积集 \(\mathbf{T}_n\)（\(n \geqslant 0\)）的并。记 \(\mathfrak{C}\) 为由紧子集生成的 \(\mathbf{T}\) 的子集族。
\end{enumerate}

\begin{enumerate}
\def\labelenumi{\alph{enumi})}
\item
  设 X 是 Lusin 空间。对于每个从 T 到 X 的 \(\mu\)-可测映射 f，存在映射序列 \(f\)\(f_m: \mathrm{T} \to \mathrm{X}\)，具有如下性质：\(\alpha\)：对于 \(\mu\)-几乎处处的 \(t \in \mathrm{T}\)，\(f(t) = \lim_{m \to \infty} f_m(t)\)；\(\beta\)：对于每个整数 m，存在将 T 划分为 \(m\)\(\mathrm{A}_j \in \mathfrak{C}\) 的有限划分，使得 \(f_m\) 在每个 \(\mathrm{A}_j\) 上为常值（化归到 \(\mathrm{T}_n\) 构成 T 的一个划分、\(\mathrm{T}_0\) 是 \(\mu\)-可忽略的、当  时 \(\mathrm{T}_n\) 是紧的且 X 是 Polish 空间的情形）。\(n \geqslant 1\)
\item
  设 \(X_{1},\ldots,X_{p}\) 是 Lusin 空间，Y 是拓扑空间，g 是从 \(X_{1}\times\cdots\times X_{p}\) 到 Y 的映射，\(f_{i}\)（\(1\leqslant i\leqslant p\)）是从 T 到 \(\mu\) 的 \(X_{i}\)-可测映射。假设对于 \(1\leqslant i\leqslant p\)，在任意固定 \(x_{i}\mapsto g(x_{1},\ldots,x_{p})\) 时，从 \(X_{i}\) 到 Y 的偏映射 \(x_{1},\ldots,x_{i-1},x_{i+1},\ldots,x_{p}\) 连续。证明从 T 到 Y 的映射 \(t\mapsto g(f_{1}(t),\ldots,f_{p}(t))\) 是 \(\mu\)-可测的。
\end{enumerate}

